\chapter{Writing the thesis} \label{ch2}

This chapter shows how the style renders the elements of a scientific text:
headings, cross-references, mathematics, theorems, tables and figures. Each
example is followed by its source code, which can be copied into the thesis.

\section{Headings and cross-references} \label{sec:headings}

The thesis is divided into chapters, sections and subsections:
\begin{Verbatim}
\chapter{Writing the thesis}  \label{ch2}
\section{Headings}            \label{sec:headings}
\subsection{Labels}
\end{Verbatim}
Subsubsections (\cs{subsubsection}) are available too, but they are neither
numbered nor listed in the table of contents: if you feel the need for
them, the section is probably too long.

Every numbered element can be given a \cs{label} and referred to with
\cs{ref} (or \cs{eqref} for equations). A non-breaking space (\verb|~|)
avoids a line break between the word and the number:
\begin{Verbatim}
as discussed in Section~\ref{sec:headings}
\end{Verbatim}
which produces ``as discussed in Section~\ref{sec:headings}''. References
are hyperlinks (Section~\ref{sec:links}).

\section{Mathematics} \label{sec:math}

Mathematics is typeset with the STIX Two Math font (see the \texttt{mathfont}
option in Section~\ref{sec:programme}). Inline formulas are
enclosed in dollar signs: \verb|$\bx_i \in \R^d$| produces
$\bx_i \in \R^d$. Displayed formulas that need a number use the
\texttt{equation} environment:
\begin{equation} \label{eq:risk}
  \min_{\bw, b} \; \frac{1}{n} \sum_{i=1}^n
    \ell\bigl(y_i, \bw^\top \bx_i + b\bigr) + \lambda \lVert \bw \rVert^2
\end{equation}
\begin{Verbatim}
\begin{equation} \label{eq:risk}
  \min_{\bw, b} \; \frac{1}{n} \sum_{i=1}^n
    \ell\bigl(y_i, \bw^\top \bx_i + b\bigr) + \lambda \lVert \bw \rVert^2
\end{equation}
\end{Verbatim}
and can be referred to as \verb|\eqref{eq:risk}|, which gives
\eqref{eq:risk}.

A displayed formula that is never referred to does not need a number:
enclose it between \verb|\[| and \verb|\]|, which centre it on a line of
its own exactly as \texttt{equation} does, but without number. For
instance, the logistic function
\[
  \sigma(z) = \frac{1}{1 + e^{-z}}
\]
is typeset by
\begin{Verbatim}
\[
  \sigma(z) = \frac{1}{1 + e^{-z}}
\]
\end{Verbatim}
The environment \texttt{equation*} is equivalent to \verb|\[|\dots\verb|\]|.
Number only the formulas you actually refer to: a thesis in which every
formula is numbered is harder to read, not easier.

Several formulas aligned on a common symbol use \texttt{align}:
\begin{align}
  \pr  &= \frac{\TP}{\TP + \FP}, &
  \rec &= \frac{\TP}{\TP + \FN}, \label{eq:prec-rec} \\
  \F_1 &= 2 \, \frac{\pr \cdot \rec}{\pr + \rec}. \label{eq:f1}
\end{align}
\begin{Verbatim}
\begin{align}
  \pr  &= \frac{\TP}{\TP + \FP}, &
  \rec &= \frac{\TP}{\TP + \FN}, \label{eq:prec-rec} \\
  \F_1 &= 2 \, \frac{\pr \cdot \rec}{\pr + \rec}. \label{eq:f1}
\end{align}
\end{Verbatim}
Its starred version \texttt{align*} aligns the formulas without numbering
any of them, while \cs{notag} at the end of a line of \texttt{align} omits
the number of that line only:
\begin{align*}
  \sigma(z)  &= \frac{1}{1 + e^{-z}}, \\
  \sigma'(z) &= \sigma(z) \bigl(1 - \sigma(z)\bigr).
\end{align*}
\begin{Verbatim}
\begin{align*}
  \sigma(z)  &= \frac{1}{1 + e^{-z}}, \\
  \sigma'(z) &= \sigma(z) \bigl(1 - \sigma(z)\bigr).
\end{align*}
\end{Verbatim}

\subsection{Notation}

Notation used throughout the thesis is best defined once, as macros in
\nolinkurl{style/symbols.tex}. The file already contains, among others, bold
vectors and matrices (\verb|\bx|, \verb|\bw|, \verb|\bX| give $\bx$, $\bw$,
$\bX$), number sets (\verb|\R| gives $\R$) and the operators used above
(\verb|\pr|, \verb|\rec|, \verb|\sgn|). Edit it freely: it belongs to your
thesis, not to the style. Note also that \cs{epsilon} is redefined as
\cs{varepsilon}, so it is printed as $\epsilon$. The meaning of the
notation can be summarised in the list of symbols
(Section~\ref{sec:symbols}).

\subsection{Theorems and definitions} \label{sec:theorems}

The style defines the theorem-like environments listed in
Table~\ref{tab:theorems}. Each numbered environment has its own counter,
which restarts at each chapter: they are numbered as chapter and
progressive number, so that the first definition and the first
proposition of Chapter~2 are Definition~2.1 and Proposition~2.1, the
second definition is Definition~2.2, and so on.

\begin{table}[htbp]
  \centering
  \caption[Theorem-like environments]{Theorem-like environments and their
    headings in the two languages.}
  \label{tab:theorems}
  \begin{tabular}{lll}
    \toprule
    Environment & English & Italian \\
    \midrule
    \texttt{thm}       & Theorem      & Teorema \\
    \texttt{lem}       & Lemma        & Lemma \\
    \texttt{prop}      & Proposition  & Proposizione \\
    \texttt{cor}       & Corollary    & Corollario \\
    \texttt{defn}, \texttt{defn*} & Definition & Definizione \\
    \texttt{fact}      & Fact         & Fatto \\
    \texttt{proof}     & Proof        & Dimostrazione \\
    \bottomrule
  \end{tabular}
\end{table}

\begin{defn} \label{def:convex}
  A function $f \colon \R^d \to \R$ is \emph{convex} if
  $f(\alpha \bx + (1 - \alpha) \by) \le \alpha f(\bx) + (1 - \alpha) f(\by)$
  for all $\bx, \by \in \R^d$ and $\alpha \in [0, 1]$.
\end{defn}

\begin{prop} \label{prop:hinge}
  For every $y \in \{-1, +1\}$, the hinge loss $\ell(y, z) = \max\{0, 1 - yz\}$
  is a convex function of $z$.
\end{prop}

\begin{proof}
  It is the maximum of two affine functions of $z$, and the maximum of convex
  functions is convex.
\end{proof}

\begin{Verbatim}
\begin{defn} \label{def:convex}
  A function $f \colon \R^d \to \R$ is \emph{convex} if ...
\end{defn}

\begin{prop} \label{prop:hinge}
  For every $y \in \{-1, +1\}$, the hinge loss ... is convex.
\end{prop}

\begin{proof}
  It is the maximum of two affine functions of $z$, ...
\end{proof}
\end{Verbatim}

\paragraph{Unnumbered definitions} Definitions that are never referred to,
for instance when recalling a standard notion, can use \texttt{defn*}: it
has the same heading as \texttt{defn}, but no number, and does not advance
the counter of definitions. Since there is no number, it cannot be the target of a
\cs{ref}.

\begin{defn*}
  The \emph{Euclidean norm} of $\bx \in \R^d$ is
  $\lVert \bx \rVert = \sqrt{x_1^2 + \dots + x_d^2}$.
\end{defn*}

\begin{Verbatim}
\begin{defn*}
  The \emph{Euclidean norm} of $\bx \in \R^d$ is
  $\lVert \bx \rVert = \sqrt{x_1^2 + \dots + x_d^2}$.
\end{defn*}
\end{Verbatim}
Unnumbered versions of the other environments can be added in the preamble
with the \cs{newtheorem*} command of the \texttt{amsthm} package; their
heading, however, does not change with the language of the programme:
\begin{Verbatim}
\newtheorem*{thm*}{Theorem}
\end{Verbatim}

\section{Tables} \label{sec:tables}

Tables are placed in a floating \texttt{table} environment, with the
caption \emph{above} the table. Use the rules of the \texttt{booktabs}
package (\cs{toprule}, \cs{midrule}, \cs{bottomrule}) and no vertical lines.
For numbers, the \texttt{S} column of \texttt{siunitx} aligns values on the
decimal point and formats uncertainties, as in Table~\ref{tab:results}.

\begin{table}[htbp]
  \centering
  \caption[Performance of the compared models]{Performance of the compared
    models on the test set (mean $\pm$ standard deviation over 10 runs). Best
    values in bold.}
  \label{tab:results}
  \begin{tabular}{l
                  S[table-format=1.3(2), detect-weight, mode=text]
                  S[table-format=1.3(2), detect-weight, mode=text]
                  S[table-format=1.3(2), detect-weight, mode=text]}
    \toprule
    {Model} & {Precision} & {Recall} & {$\F_1$} \\
    \midrule
    Logistic regression & 0.812(14) & 0.774(21) & 0.792(12) \\
    SVM (linear)        & 0.828(11) & 0.781(18) & 0.803(10) \\
    SVM (RBF)           & 0.851(09) & 0.802(15) & 0.825(08) \\
    Random forest       & \bfseries 0.869(10) & \bfseries 0.815(13) & \bfseries 0.841(07) \\
    \bottomrule
  \end{tabular}
\end{table}

\begin{Verbatim}
\begin{table}[htbp]
  \centering
  \caption{Performance of the compared models ...}
  \label{tab:results}
  \begin{tabular}{l
                  S[table-format=1.3(2), detect-weight, mode=text]
                  ...}
    \toprule
    {Model} & {Precision} & {Recall} & {$\F_1$} \\
    \midrule
    Logistic regression & 0.812(14) & 0.774(21) & 0.792(12) \\
    ...
    Random forest       & \bfseries 0.869(10) & ... \\
    \bottomrule
  \end{tabular}
\end{table}
\end{Verbatim}
Headers in \texttt{S} columns go between braces; the options
\texttt{detect-weight} and \texttt{mode=text} are needed to typeset a value
in bold with \cs{bfseries}.

\subsection{Rotated tables} \label{sec:rotated}

A table too wide for the page can be placed on a page in landscape
orientation, so that its width runs along the height of the sheet. The
\texttt{landscape} environment of the \texttt{pdflscape} package, loaded in
\texttt{main.tex}, typesets its content on rotated pages; in the PDF these
pages are also marked as landscape, so that viewers display them the right
way up. Wrapping it in \cs{afterpage} defers the landscape page until the
end of the current one, so that the text flows without leaving a gap, as for
Table~\ref{tab:datasets}.

\afterpage{%
\begin{landscape}
\begin{table}
  \centering
  \caption[Performance on three datasets]{Precision, recall and $\F_1$ score
    of the compared models on three datasets (mean over 10 runs). Best values
    for each dataset in bold.}
  \label{tab:datasets}
  \sisetup{table-format=1.3, detect-weight, mode=text}
  \begin{tabular}{l SSS SSS SSS}
    \toprule
    & \multicolumn{3}{c}{Breast cancer}
    & \multicolumn{3}{c}{Credit default}
    & \multicolumn{3}{c}{Hospital readmission} \\
    \cmidrule(lr){2-4} \cmidrule(lr){5-7} \cmidrule(lr){8-10}
    {Model} & {Precision} & {Recall} & {$\F_1$}
            & {Precision} & {Recall} & {$\F_1$}
            & {Precision} & {Recall} & {$\F_1$} \\
    \midrule
    Logistic regression & 0.951 & 0.962 & 0.956
                        & 0.692 & 0.358 & 0.472
                        & 0.618 & 0.540 & 0.576 \\
    SVM (linear)        & 0.958 & 0.965 & 0.961
                        & 0.701 & 0.341 & 0.459
                        & 0.622 & 0.533 & 0.574 \\
    SVM (RBF)           & \bfseries 0.972 & 0.969 & \bfseries 0.970
                        & 0.688 & 0.372 & 0.483
                        & 0.635 & 0.548 & 0.588 \\
    Random forest       & 0.965 & \bfseries 0.971 & 0.968
                        & \bfseries 0.712 & \bfseries 0.389 & \bfseries 0.503
                        & \bfseries 0.651 & \bfseries 0.566 & \bfseries 0.606 \\
    \bottomrule
  \end{tabular}
\end{table}
\end{landscape}}

\begin{Verbatim}
\afterpage{%
\begin{landscape}
\begin{table}
  \centering
  \caption{Precision, recall and $\F_1$ score of the compared models ...}
  \label{tab:datasets}
  \sisetup{table-format=1.3, detect-weight, mode=text}
  \begin{tabular}{l SSS SSS SSS}
    \toprule
    & \multicolumn{3}{c}{Breast cancer}
    & \multicolumn{3}{c}{Credit default}
    & \multicolumn{3}{c}{Hospital readmission} \\
    \cmidrule(lr){2-4} \cmidrule(lr){5-7} \cmidrule(lr){8-10}
    {Model} & {Precision} & {Recall} & {$\F_1$} & ... \\
    \midrule
    Logistic regression & 0.951 & 0.962 & 0.956 & ... \\
    ...
    \bottomrule
  \end{tabular}
\end{table}
\end{landscape}}
\end{Verbatim}
Here \cs{sisetup} sets the options of all the \texttt{S} columns at once,
\cs{multicolumn} spans a header over several columns and \cs{cmidrule} draws a
partial rule below it. Wide figures are placed on landscape pages in the
same way, with a \texttt{figure} environment inside \texttt{landscape}.

\subsection{Tables spanning several pages} \label{sec:longtable}

A table longer than a page is written with the \texttt{longtable}
environment of the package with the same name, loaded in
\texttt{main.tex}. Unlike \texttt{table}, it does not float: it starts where
it is written and continues on as many pages as needed. Its first rows
declare what is repeated on each page: the header of the first page (ended
by \cs{endfirsthead}), the header of the following pages (\cs{endhead}),
the footer of all pages but the last (\cs{endfoot}) and the footer of the
last page (\cs{endlastfoot}). Table~\ref{tab:grid} shows the result, and its
source follows; the values explored by the search are listed in
Table~\ref{tab:search-space}, in Appendix~\ref{app:appendices}.

\begin{longtable}{r r r l S[table-format=1.3, detect-weight, mode=text]}
  \caption[Grid search of a random forest]{$\F_1$ score of a random forest
    for each configuration of a grid search over its hyperparameters
    (illustrative values). Best value in bold.}
  \label{tab:grid} \\
  \toprule
  {\#} & {Max depth} & {Trees} & {Max features} & {$\F_1$} \\
  \midrule
  \endfirsthead
  \caption[]{(continued)} \\
  \toprule
  {\#} & {Max depth} & {Trees} & {Max features} & {$\F_1$} \\
  \midrule
  \endhead
  \midrule
  \multicolumn{5}{r}{\emph{Continued on next page}} \\
  \endfoot
  \bottomrule
  \endlastfoot
    1 & 4 & 50 & \texttt{sqrt} & 0.775 \\
    2 & 4 & 50 & \texttt{log2} & 0.769 \\
    3 & 4 & 50 & 0.5 & 0.772 \\
    4 & 4 & 100 & \texttt{sqrt} & 0.777 \\
    5 & 4 & 100 & \texttt{log2} & 0.771 \\
    6 & 4 & 100 & 0.5 & 0.780 \\
    7 & 4 & 200 & \texttt{sqrt} & 0.783 \\
    8 & 4 & 200 & \texttt{log2} & 0.777 \\
    9 & 4 & 200 & 0.5 & 0.779 \\
    10 & 4 & 500 & \texttt{sqrt} & 0.787 \\
    11 & 4 & 500 & \texttt{log2} & 0.781 \\
    12 & 4 & 500 & 0.5 & 0.784 \\
    13 & 8 & 50 & \texttt{sqrt} & 0.800 \\
    14 & 8 & 50 & \texttt{log2} & 0.794 \\
    15 & 8 & 50 & 0.5 & 0.797 \\
    16 & 8 & 100 & \texttt{sqrt} & 0.802 \\
    17 & 8 & 100 & \texttt{log2} & 0.796 \\
    18 & 8 & 100 & 0.5 & 0.805 \\
    19 & 8 & 200 & \texttt{sqrt} & 0.808 \\
    20 & 8 & 200 & \texttt{log2} & 0.802 \\
    21 & 8 & 200 & 0.5 & 0.804 \\
    22 & 8 & 500 & \texttt{sqrt} & 0.812 \\
    23 & 8 & 500 & \texttt{log2} & 0.806 \\
    24 & 8 & 500 & 0.5 & 0.809 \\
    25 & 16 & 50 & \texttt{sqrt} & 0.809 \\
    26 & 16 & 50 & \texttt{log2} & 0.803 \\
    27 & 16 & 50 & 0.5 & 0.812 \\
    28 & 16 & 100 & \texttt{sqrt} & 0.817 \\
    29 & 16 & 100 & \texttt{log2} & 0.811 \\
    30 & 16 & 100 & 0.5 & 0.813 \\
    31 & 16 & 200 & \texttt{sqrt} & 0.816 \\
    32 & 16 & 200 & \texttt{log2} & 0.810 \\
    33 & 16 & 200 & 0.5 & 0.819 \\
    34 & 16 & 500 & \texttt{sqrt} & 0.821 \\
    35 & 16 & 500 & \texttt{log2} & 0.815 \\
    36 & 16 & 500 & 0.5 & \bfseries 0.824 \\
    37 & 32 & 50 & \texttt{sqrt} & 0.807 \\
    38 & 32 & 50 & \texttt{log2} & 0.801 \\
    39 & 32 & 50 & 0.5 & 0.810 \\
    40 & 32 & 100 & \texttt{sqrt} & 0.815 \\
    41 & 32 & 100 & \texttt{log2} & 0.809 \\
    42 & 32 & 100 & 0.5 & 0.811 \\
    43 & 32 & 200 & \texttt{sqrt} & 0.814 \\
    44 & 32 & 200 & \texttt{log2} & 0.808 \\
    45 & 32 & 200 & 0.5 & 0.817 \\
    46 & 32 & 500 & \texttt{sqrt} & 0.819 \\
    47 & 32 & 500 & \texttt{log2} & 0.813 \\
    48 & 32 & 500 & 0.5 & 0.822 \\
    49 & \texttt{None} & 50 & \texttt{sqrt} & 0.807 \\
    50 & \texttt{None} & 50 & \texttt{log2} & 0.801 \\
    51 & \texttt{None} & 50 & 0.5 & 0.804 \\
    52 & \texttt{None} & 100 & \texttt{sqrt} & 0.809 \\
    53 & \texttt{None} & 100 & \texttt{log2} & 0.803 \\
    54 & \texttt{None} & 100 & 0.5 & 0.806 \\
    55 & \texttt{None} & 200 & \texttt{sqrt} & 0.815 \\
    56 & \texttt{None} & 200 & \texttt{log2} & 0.809 \\
    57 & \texttt{None} & 200 & 0.5 & 0.811 \\
    58 & \texttt{None} & 500 & \texttt{sqrt} & 0.819 \\
    59 & \texttt{None} & 500 & \texttt{log2} & 0.813 \\
    60 & \texttt{None} & 500 & 0.5 & 0.816 \\
\end{longtable}

\begin{Verbatim}
\begin{longtable}{r r r l S[table-format=1.3, detect-weight, mode=text]}
  \caption{$\F_1$ score of a random forest for each configuration ...}
  \label{tab:grid} \\
  \toprule
  {\#} & {Max depth} & {Trees} & {Max features} & {$\F_1$} \\
  \midrule
  \endfirsthead
  \caption[]{(continued)} \\
  \toprule
  {\#} & {Max depth} & {Trees} & {Max features} & {$\F_1$} \\
  \midrule
  \endhead
  \midrule
  \multicolumn{5}{r}{\emph{Continued on next page}} \\
  \endfoot
  \bottomrule
  \endlastfoot
  1 & 4 & 50 & \texttt{sqrt} & 0.775 \\
  2 & 4 & 50 & \texttt{log2} & 0.769 \\
  ...
\end{longtable}
\end{Verbatim}
The caption and the label are the first row of the table, ended by
\verb|\\|. On the following pages, \verb|\caption[]{(continued)}| repeats
the table number followed by ``(continued)'', without adding a second entry
to the list of tables. Since a long table cannot float, introduce it at the
point of the text where it belongs, and keep it out of paragraphs: leave an
empty line before and after it.

\section{Figures} \label{sec:figures}

Figures are placed in a floating \texttt{figure} environment, with the
caption \emph{below} the figure. Images are included with
\cs{includegraphics}, looking for them in the \texttt{images/} directory;
vector formats (PDF) are preferable to bitmaps:
\begin{Verbatim}
\begin{figure}[htbp]
  \centering
  \includegraphics[width=0.8\textwidth]{kde_plot.pdf}
  \caption{Kernel density estimate of ...}
  \label{fig:kde}
\end{figure}
\end{Verbatim}

Plots can also be drawn directly in \LaTeX{} with \texttt{pgfplots}, so
that their fonts match the text, as in Figure~\ref{fig:learning-curve}.
The colours defined in \nolinkurl{style/colors.tex}, as well as
\texttt{slimblue}, can be used in plots and diagrams.

\begin{figure}[htbp]
  \centering
  \begin{tikzpicture}
    \begin{axis}[
        width=0.8\textwidth, height=0.5\textwidth,
        xlabel={Training set size}, ylabel={$\F_1$},
        xmin=0, xmax=10000, ymin=0.5, ymax=0.9,
        scaled x ticks=false,
        xticklabel style={/pgf/number format/1000 sep={\,}},
        grid=major, grid style={gray!20},
        legend pos=south east, legend cell align=left,
        every axis plot/.append style={thick, mark size=1.5pt},
      ]
      \addplot[color=pastelBlueDark, mark=*]
        coordinates {(500,0.58) (1000,0.65) (2000,0.72) (4000,0.77)
                     (6000,0.79) (8000,0.80) (10000,0.80)};
      \addlegendentry{Logistic regression}
      \addplot[color=attnOrangeDark, mark=square*]
        coordinates {(500,0.55) (1000,0.63) (2000,0.73) (4000,0.80)
                     (6000,0.82) (8000,0.83) (10000,0.84)};
      \addlegendentry{Random forest}
    \end{axis}
  \end{tikzpicture}
  \caption[Learning curves]{Learning curves: $\F_1$ score on the test set as
    a function of the training set size.}
  \label{fig:learning-curve}
\end{figure}

\begin{Verbatim}
\begin{figure}[htbp]
  \centering
  \begin{tikzpicture}
    \begin{axis}[width=0.8\textwidth, height=0.5\textwidth,
                 xlabel={Training set size}, ylabel={$\F_1$}, ...]
      \addplot[color=pastelBlueDark, mark=*]
        coordinates {(500,0.58) (1000,0.65) ... (10000,0.80)};
      \addlegendentry{Logistic regression}
      ...
    \end{axis}
  \end{tikzpicture}
  \caption{Learning curves: ...}
  \label{fig:learning-curve}
\end{figure}
\end{Verbatim}

\subsection{Subfigures} \label{sec:subfigures}

Related images can be grouped in a single figure, each with its own
caption, using the \texttt{subfigure} environment of the
\texttt{subcaption} package. Each subfigure is given a width, here slightly
less than half of the line so that two of them fit side by side, and
everything inside it, including plots, should be sized relative to
\cs{linewidth}. Subfigures are labelled (a), (b), and so on: a reference to
a subfigure, such as \verb|\ref{fig:sigmoid}|, gives
\ref{fig:sigmoid}, while \verb|\subref{fig:sigmoid}| gives just the label,
\subref{fig:sigmoid}, as in Figure~\ref{fig:subfigures}.

\begin{figure}[htbp]
  \centering
  \begin{subfigure}[b]{0.48\textwidth}
    \centering
    \begin{tikzpicture}
      \begin{axis}[
          width=0.95\linewidth, height=0.8\linewidth,
          xlabel={$z$}, ylabel={$\sigma(z)$},
          xmin=-6, xmax=6, ymin=0, ymax=1,
          grid=major, grid style={gray!20},
        ]
        \addplot[color=pastelBlueDark, thick, domain=-6:6, samples=100]
          {1/(1+exp(-x))};
      \end{axis}
    \end{tikzpicture}
    \caption{Logistic function.}
    \label{fig:sigmoid}
  \end{subfigure}
  \hfill
  \begin{subfigure}[b]{0.48\textwidth}
    \centering
    \begin{tikzpicture}
      \begin{axis}[
          width=0.95\linewidth, height=0.8\linewidth,
          xlabel={False positive rate}, ylabel={True positive rate},
          xmin=0, xmax=1, ymin=0, ymax=1,
          grid=major, grid style={gray!20},
        ]
        \addplot[color=gray, dashed, domain=0:1] {x};
        \addplot[color=attnOrangeDark, thick, mark=*, mark size=1.2pt]
          coordinates {(0,0) (0.02,0.35) (0.05,0.58) (0.10,0.74) (0.20,0.87)
                       (0.35,0.94) (0.60,0.98) (1,1)};
      \end{axis}
    \end{tikzpicture}
    \caption{ROC curve of the random forest.}
    \label{fig:roc}
  \end{subfigure}
  \caption[Logistic function and ROC curve]{Two subfigures:
    \subref{fig:sigmoid} the logistic function and \subref{fig:roc} a ROC
    curve, with the diagonal of a random classifier shown dashed.}
  \label{fig:subfigures}
\end{figure}

\begin{Verbatim}
\begin{figure}[htbp]
  \centering
  \begin{subfigure}[b]{0.48\textwidth}
    \centering
    \begin{tikzpicture}
      \begin{axis}[width=0.95\linewidth, height=0.8\linewidth, ...]
        \addplot[color=pastelBlueDark, thick, domain=-6:6]
          {1/(1+exp(-x))};
      \end{axis}
    \end{tikzpicture}
    \caption{Logistic function.}
    \label{fig:sigmoid}
  \end{subfigure}
  \hfill
  \begin{subfigure}[b]{0.48\textwidth}
    ...
    \caption{ROC curve of the random forest.}
    \label{fig:roc}
  \end{subfigure}
  \caption{Two subfigures: \subref{fig:sigmoid} the logistic function
    and \subref{fig:roc} a ROC curve, ...}
  \label{fig:subfigures}
\end{figure}
\end{Verbatim}
The optional argument \texttt{[b]} aligns the subfigures at their bottom,
so that their captions are on the same line even if the images have
different heights. Images are included in the same way, with
\verb|\includegraphics[width=\linewidth]{...}| inside each subfigure.

\section{Algorithms} \label{sec:algorithms}

Algorithms are written in pseudocode with the \texttt{algpseudocode}
package, inside the floating \texttt{algorithm} environment of the
\texttt{algorithm} package, both loaded in \texttt{main.tex}. Like figures
and tables, algorithms have a caption and are numbered within each
chapter, as Algorithm~\ref{alg:kmeans}.

\begin{algorithm}[htbp]
  \caption{$k$-means clustering}
  \label{alg:kmeans}
  \begin{algorithmic}[1]
    \Require points $\bx_1, \dots, \bx_n \in \R^d$, number of clusters $k$
    \Ensure centroids $\symbfit{\mu}_1, \dots, \symbfit{\mu}_k$ and
      assignments $c_1, \dots, c_n$
    \State choose $\symbfit{\mu}_1, \dots, \symbfit{\mu}_k$ at random among
      the points
    \Repeat
      \For{$i = 1, \dots, n$}
        \State $c_i \gets \operatorname*{arg\,min}_j
          \lVert \bx_i - \symbfit{\mu}_j \rVert$
          \Comment{nearest centroid}
      \EndFor
      \For{$j = 1, \dots, k$}
        \State $\symbfit{\mu}_j \gets$ mean of the points $\bx_i$ with
          $c_i = j$ \label{alg:kmeans:update}
      \EndFor
    \Until{no assignment $c_i$ changes}
    \State \Return $\symbfit{\mu}_1, \dots, \symbfit{\mu}_k$,
      $c_1, \dots, c_n$
  \end{algorithmic}
\end{algorithm}

\begin{Verbatim}
\begin{algorithm}[htbp]
  \caption{$k$-means clustering}
  \label{alg:kmeans}
  \begin{algorithmic}[1]
    \Require points $\bx_1, \dots, \bx_n \in \R^d$, number of clusters $k$
    \Ensure centroids ... and assignments $c_1, \dots, c_n$
    \State choose $\symbfit{\mu}_1, \dots, \symbfit{\mu}_k$ at random ...
    \Repeat
      \For{$i = 1, \dots, n$}
        \State $c_i \gets \operatorname*{arg\,min}_j
          \lVert \bx_i - \symbfit{\mu}_j \rVert$
          \Comment{nearest centroid}
      \EndFor
      \For{$j = 1, \dots, k$}
        \State $\symbfit{\mu}_j \gets$ mean of the points ...
          \label{alg:kmeans:update}
      \EndFor
    \Until{no assignment $c_i$ changes}
    \State \Return $\symbfit{\mu}_1, \dots, \symbfit{\mu}_k$,
      $c_1, \dots, c_n$
  \end{algorithmic}
\end{algorithm}
\end{Verbatim}
The optional argument \texttt{[1]} of \texttt{algorithmic} numbers every
line; a line can then be given a \cs{label} and referred to, as
line~\ref{alg:kmeans:update} of Algorithm~\ref{alg:kmeans}, which updates
the centroids. \cs{Require} and \cs{Ensure} introduce the input and the
output (printed as ``Input:'' and ``Output:''), \cs{State} a generic
step, \cs{Comment} a comment aligned to the right; besides \cs{For} and \cs{Repeat}--\cs{Until}, the package provides
\cs{If}--\cs{ElsIf}--\cs{Else}--\cs{EndIf}, \cs{While}--\cs{EndWhile},
\cs{Function}--\cs{EndFunction} and \cs{Call}. The heading is
``Algorithm'' or ``Algoritmo'' according to the language. Algorithms are
collected in
the list of algorithms, printed by \cs{listofalgorithms} after the list of
tables (Section~\ref{sec:structure}).
